% =============================================================================
% Project : THE PHILOSOPHY OF VIABILITY
% Book    : Book I — Ontology of Viability
% Book ID : VS.B01
% Part    : Part I — Identity and Manifestation
% Part ID : VS.B01.P01
% Chapter : Chapter 08 — Comparative Philosophical Examination
% Chapter ID : VS.B01.P01.C08
%
% File :
%   4_ASSETS/tables/
%   VS.B01.P01.C08.S15.SS18.P004.TL01.averroes_claims_and_comparative_limits.tex
%
% Asset :
%   TL01 — Averroes's Principal Claims and Comparative Limits
%
% Structural Identifier :
%   VS.B01.P01.C08.S15.SS18.P004.TL01
%
% Purpose :
%   To distinguish the functional or logical form of the
%   reconstructed claims from their epistemic or justificatory
%   status, and to record their limits for Book I.
%
% Dependencies :
%   longtable, booktabs, array.
%   These packages are already loaded by the Book I preamble.
%
% =============================================================================

\begingroup
\small
\renewcommand{\arraystretch}{1.18}
\setlength{\tabcolsep}{4pt}
\setlength{\LTcapwidth}{\linewidth}

\def\colw#1{\dimexpr(\linewidth-8\tabcolsep)*#1/100\relax}

\def\rowsep{\specialrule{0.2pt}{3pt}{3pt}}

\def\tablehead{%
    \toprule
    \textbf{Claim or issue} &
    \textbf{Averroist reconstruction} &
    \textbf{Functional or logical form} &
    \textbf{Epistemic or justificatory status} &
    \textbf{Comparative limit for Book I}
    \\
    \midrule
}

\def\grouprow#1{%
    \multicolumn{5}{@{}l}{\footnotesize\textsc{#1}} \\
    \cmidrule{1-5}
    \addlinespace[1pt]
}

\begin{longtable}{@{}
    >{\raggedright\arraybackslash}p{\colw{14}}
    >{\raggedright\arraybackslash}p{\colw{26}}
    >{\raggedright\arraybackslash}p{\colw{14}}
    >{\raggedright\arraybackslash}p{\colw{17}}
    >{\raggedright\arraybackslash}p{\colw{29}}
@{}}

\caption{Averroes's Principal Claims and Their Comparative Limits}
\label{tab:VS.B01.P01.C08.S15.SS18.P004.TL01}
\\

\tablehead
\endfirsthead

\multicolumn{5}{@{}l}{%
    \footnotesize\textit{
        Table~\thetable\ (continued):
        Averroes's Principal Claims and Their Comparative Limits
    }
}
\\
\tablehead
\endhead

\multicolumn{5}{@{}r}{%
    \footnotesize\textit{Continued on next page}
}
\\
\endfoot

\bottomrule
\endlastfoot

\grouprow{Being and substance}

\textbf{Being qua being}
&
Metaphysics investigates being universally and the principles
that belong to existents across the special sciences.
&
Thesis
&
Argued/Derived Claims
&
The subject of metaphysics does not by itself settle
the relation between Being and Identity in Book I.
\\
\rowsep

\textbf{Substance}
&
Substance occupies a primary position among the modes
of being; accidents depend upon substance.
&
Thesis
&
Argued/Derived Claims
&
Substance is not thereby identical with Identity.
\\
\rowsep

\grouprow{Matter, form, and essence}

\textbf{Prime matter}
&
Matter is considered as a potential subject for
substantial forms and natural change.
&
Thesis
&
Argued/Derived Claims
&
Material potentiality is not Identity or Viability.
\\
\rowsep

\textbf{Substantial form}
&
Form actualizes matter and accounts for the determinate
nature of the composite substance.
&
Thesis
&
Argued/Derived Claims
&
Formal determination is not equivalent to Identity.
\\
\rowsep

\textbf{Essence and quiddity}
&
The essence of a sensible substance is explained
through its substantial form.
&
Proposition
&
Conditionally Supported Claims
&
Essence must not be renamed Identity.
\\
\rowsep

\grouprow{Existence and Avicenna}

\textbf{Existence and essence}
&
Averroes criticizes the treatment of existence as
a common attribute or accident of essence.
&
Thesis
&
Conditionally Supported Claims
&
The interpretation of Avicenna's position remains
disputed; the critique is not treated as decisive.
\\
\rowsep

\textbf{Avicennian distinction}
&
Avicenna distinguishes quiddity considered in itself
from the actuality of existence and its causal conditions.
&
Thesis
&
Argued/Derived Claims
&
The distinction is not equivalent to the Being--Identity
distinction.
\\
\rowsep

\grouprow{Unity, individuation, and change}

\textbf{Unity and multiplicity}
&
The one and the many are examined as central
metaphysical problems within the Aristotelian framework.
&
Thesis
&
Argued/Derived Claims
&
Unity is not identical with Identity, Singularity,
or Aggregate.
\\
\rowsep

\textbf{Material multiplication}
&
Extended matter contributes to the multiplication
of material forms.
&
Proposition
&
Conditionally Supported Claims
&
Multiplication is not a universal principle of Identity.
\\
\rowsep

\textbf{Individuation}
&
Forms are individual in their existence; matter
contributes to material plurality.
&
Thesis
&
Conditionally Supported Claims
&
Individuation is not equivalent to Identity.
\\
\rowsep

\textbf{Accidental change}
&
A substance may gain or lose accidents while remaining
the subject of those changes.
&
Proposition
&
Argued/Derived Claims
&
Persistence through accidents is not a definition of Identity.
\\
\rowsep

\textbf{Substantial change}
&
The departure of substantial form marks the end
of the former substance or its continuation only
homonymously.
&
Thesis
&
Conditionally Supported Claims
&
No universal Book I criterion of Identity follows
without further argument.
\\
\rowsep

\grouprow{Causality and cosmology}

\textbf{Natural causality}
&
Natural agents actualize material potentialities
through causal processes.
&
Thesis
&
Argued/Derived Claims
&
Causal dependence does not establish identity with
the cause or the causal process.
\\
\rowsep

\textbf{Eternity of the world}
&
Averroes defends an Aristotelian account of eternal
motion and the world.
&
Thesis
&
Conditionally Supported Claims
&
This is a historical cosmological commitment,
not a foundational commitment of Book I.
\\
\rowsep

\grouprow{Comparative status}

\textbf{The so-called two truths}
&
The inherited description of Averroes as affirming
two contradictory truths is rejected.
&
Rejected Thesis
&
Rejected/Superseded Claims
&
The label is not used as a characterization of
Averroes's position.
\\
\rowsep

\textbf{Being and Identity}
&
Book I maintains Analytical Distinction, Ontological
Inseparability, and Absence of Ontological Priority.
&
Foundational Commitment
&
Foundationally Adopted Commitments
&
These three commitments remain separate and are
not derived from Averroes.
\\
\rowsep

\textbf{Positive role of Identity}
&
The explanatory role of Identity beyond substantial
form, essence, individuation, and persistence
requires further development.
&
Open Problem
&
Unresolved/Open Problems
&
The positive content of Identity must be established
rather than assumed.
\\
\rowsep

\textbf{Comparative synthesis}
&
Averroes's ontology and Book I address overlapping
problems through distinct explanatory architectures.
&
Thesis
&
Provisional Constructions
&
The comparison establishes critical proximity,
not conceptual equivalence.
\\

\end{longtable}

\endgroup

% =============================================================================
% END OF TABLE TL01
% =============================================================================